% SGA 3, Expose VII_B, English reconstruction -- Loop 1.
% Scope: Section 5.2.3 in full on authority local pages 92--93;
% stops immediately before Section 5.2.4.
% Combined-reader pages: 605--606. Re-edition printed pages: 594--595.
% Original-edition marginal locator in scope: 560.
% Editorial notes in scope: 158--159.
% Authority: Polo--Gille Exp7B-13oct24.pdf,
% SHA-256 D303690E0195654BBE7C3DAA5D89EA23E24CD709A3575217234AC19AE8836694.
% Jacob Reinhold's CC BY 4.0 Markdown was consulted as comparison lineage.
% Three genuine diagrams are preserved as two exact PDF-derived PNGs
% for Loop 1 and remain due for native reconstruction in Loop 2.

\subsubsection*{5.2.3}
\phantomsection
\label{sga3:VIIB:5.2.3}

现在假设 \(B\) 的高度有限且至多为 \(r\)。令
\[
  \pi:k\xrightarrow{\ \sim\ }k,
  \qquad
  x\longmapsto x^p,
\]
为 Frobenius 自同构。线性映射
\[
  B\widehat\otimes_{\pi}k\longrightarrow B,
  \qquad
  b\widehat\otimes_{\pi}x
  \longmapsto
  b^px=(bx^{1/p})^p,
\]
的像闭，等于 \(B\) 的闭子代数
\[
  B^p=\{\,b^p\mid b\in B\,\}
\]
置
\[
  J=B^p\cap I=B^p\cap I_A.
\]

\footnote{\label{sga3:VIIB:nde:158}编者增加了以下内容。}
令 \(G_1\) 为 Frobenius 态射
\[
  \operatorname{Fr}:G\longrightarrow G^{(p)},
\]
的核，令 \(HG_1\) 为形式子群 \(H^{(p)}\subset G^{(p)}\) 在 \(G\) 中的逆像形式子群。则 \(HG_1\) 由 \(AI\) 的元素的 \(p\) 次幂生成的闭理想定义；此理想为 \(AJ\)。此外，\(G/H\) 的构造与基变换交换，因为 \(G/H\) 表示平坦拓扑下的商层\textup{（参见 \SGARefLink{sga3:VIIB:2.4}{2.4}）}。因此，
\[
  (G/H)^{(p)}=G^{(p)}/H^{(p)},
\]
并有如下交换图：

\begin{figure}[H]
  \centering
  % SGA3-VII-DIAG-132
  % SGA3-VII-DIAG-132
% Authority: Exp7B-13oct24.pdf, local p. 92.
% Crop: figures/Exp7B_localp092_frobenius_quotient_dual_diagrams.png
% Locator: tex/components/88_expose_VIIB_section523_finite_height_induction_en.tex, Loop 1 line 60.
\begin{tikzcd}[row sep=2.5em,column sep=2.8em]
G \arrow[r,"\operatorname{Fr}"] \arrow[d]
  & G^{(p)} \arrow[d] \\
G/H \arrow[r,"\operatorname{Fr}"']
  & G^{(p)}/H^{(p)}
\end{tikzcd}
\qquad
\begin{tikzcd}[row sep=2.5em,column sep=2.8em]
A
  & A\widehat\otimes_\pi k
      \arrow[l,"{a\widehat\otimes1\mapsto a^p}"'] \\
B \arrow[u]
  & B\widehat\otimes_\pi k
      \arrow[l,"{b\widehat\otimes1\mapsto b^p}"']
      \arrow[u]
\end{tikzcd}

\end{figure}

由此可知 \(B^p\) 是商 \(G/HG_1\) 的仿射代数。\footnote{\label{sga3:VIIB:nde:159}如原文指出的，这也可由 \SGARefLink{sga3:VIIB:5.1:proposition}{命题~5.1} 推出。编者更愿意给出上述论证，它不用该命题的蕴含关系 \textup{\SGARefLink{sga3:vii:paragraph:5-1:equation:i}{(i)}} \(\Rightarrow\) \textup{\SGARefLink{sga3:VIIB:5.1:proposition:item:ii}{(ii)}}。}
暂用 \(C\) 表示商 \(HG_1/H\) 的仿射代数。由于 \(G/H\) 的构造与基变换交换，方图

\begin{figure}[H]
  \centering
  % SGA3-VII-DIAG-133
  % SGA3-VII-DIAG-133
% Authority: Exp7B-13oct24.pdf, local p. 92.
% Crop: figures/Exp7B_localp092_cartesian_HG1_square.png
% Locator: tex/components/88_expose_VIIB_section523_finite_height_induction_en.tex, Loop 1 line 76.
\begin{tikzcd}[row sep=2.8em,column sep=3.5em]
HG_1 \arrow[r] \arrow[d]
  & G \arrow[d] \\
HG_1/H \arrow[r]
  & G/H
\end{tikzcd}

\end{figure}

\noindent
为笛卡尔图。从而
\[
  A\widehat\otimes_BC
  \simeq
  A/AJ
  =
  A\widehat\otimes_B(B/BJ).
\]
由于 \(A\) 在 \(B\) 上拓扑自由\textup{（由 2.4，因为 \(A\) 和 \(B\) 为局部的）}，自然态射
\[
  B/BJ\longrightarrow C
\]
为同构。因此 \(B/BJ\) 是 \(HG_1/H\) 的仿射代数，并且显然
\[
  B/BJ=B_1
\]
的高度至多为 \(1\)，因为
\[
  J=\bigl((x^p)\bigr)_{x\in\mathfrak r(B)}.
\]

令
\[
  B'
  =
  k[[\omega]]
  \big/
  \bigl((x^{p^r})\bigr)_{x\in\omega_r},
\]
令 \(\phi:B'\to B\) 为 \SGARefLink{sga3:VIIB:5.2.1:section-proposition}{5.2.1} 中引入的态射，令
\[
  B'^p=\{\,x^p\mid x\in B'\,\},
\]
并令 \(J'\) 为 \(B'^p\) 的增广理想。于是有交换图

\begin{figure}[H]
  \centering
  % SGA3-VII-DIAG-134
  % SGA3-VII-DIAG-134
% Authority: Exp7B-13oct24.pdf, local p. 93.
% Crop: figures/Exp7B_localp093_phi_reduction_square.png
% Locator: tex/components/88_expose_VIIB_section523_finite_height_induction_en.tex, Loop 1 line 122.
\begin{tikzcd}[row sep=2.8em,column sep=3.7em]
B' \arrow[r,"\psi"] \arrow[d]
  & B \arrow[d] \\
B'_1=B'/\overline{B'J'}
  \arrow[r,"\phi_1"]
  & B_1=B/\overline{BJ}
\end{tikzcd}

\end{figure}

\noindent
且由 \SGARefLink{sga3:VIIB:5.2.2}{5.2.2}，\(\phi_1\) 为同构。

另一方面，由 2.4，\(A\) 在
\[
  B=\mathscr A(G/H)
  \qquad\text{并在}\qquad
  B^p=\mathscr A(G/HG_1).
\]
上拓扑平坦。由 \SGARefLink{sga3:VIIB:1.3.3}{1.3.3}，可知 \(B\) 在 \(B^p\) 上拓扑平坦。此外，下文 \SGARefLink{sga3:VIIB:5.2.4}{5.2.4} 表明，\(\phi\) 诱导的态射
\[
  B'^p\longrightarrow B^p
\]
为同构。因此可将 \SGARefLink{sga3:VIIB:0.3.4}{0.3.4} 用于伪紧环 \(B'^p=B^p\) 和伪紧 \(B^p\)-模 \(M=B'\)、\(N=B\)。由上所述，
\[
  \phi_1
  =
  \phi\widehat\otimes_{B^p}k
\]
为同构，从而 \(\phi\) 本身为同构。结合下文 \SGARefLink{sga3:VIIB:5.2.4}{5.2.4}，这就证明了 \(B\) 高度有限时的 \SGARefLink{sga3:VIIB:5.2:theorem}{5.2}。
